\documentclass[12pt]{article}
\usepackage{amsmath,amssymb}

\begin{document}
\section*{{2024 AIME II Problems/Problem 12}}
Source: AoPS Wiki

\subsection*{Solution}
Begin by finding the equation of the line $\overline{AB}$ : $y= -\sqrt{3}x+\frac{\sqrt{3}}{2}$ Now, consider the general equation of all lines that belong to $\mathcal{F}$ . Let $P$ be located at $(a,0)$ and $Q$ be located at $(0,b)$ . With these assumptions, we may arrive at the equation $ay +bx =ab$ . However, a critical condition that must be satisfied by our parameters is that $a^2+b^2=1$ , since the length of $\overline{PQ}=1$ . Here's the golden trick that resolves the problem: we wish to find some point $C$ along $\overline{AB}$ such that $\overline{PQ}$ passes through $C$ if and only if $a=\frac{1}{2}$ . It's not hard to convince oneself of this, since the property $a^2+b^2=1$ implies that if $a=\frac{1}{2}$ , then $\overline{PQ}=\overline{AB}$ . We should now try to relate the point $C$ to some value of $a$ . This is accomplished by finding the intersection of two lines: \[ a\left(-\sqrt{3}x +\frac{\sqrt{3}}{2}\right) + x\sqrt{1-a^2} = a\sqrt{1-a^2} \] Where we have also used the fact that $b=\sqrt{1-a^2}$ , which follows nicely from $a^2+b^2 =1$ . \[ a\left(-\sqrt{3}x +\frac{\sqrt{3}}{2}\right) = (a-x)\sqrt{1-a^2} \] Square both sides and go through some algebraic manipulations to arrive at \[ -a^4 +2xa^3+\left(-4x^2+3x+\frac{1}{4}\right)a^2-2xa+x^2=0 \] Note how $a=\frac{1}{2}$ is a solution to this polynomial, and it is logically so. If we found the set of intersections consisting of line segment $\overline{AB}$ with an identical copy of itself, every single point on the line (all $x$ values) should satisfy the equation. Thus, we can perform polynomial division to eliminate the extraneous solution $a=\frac{1}{2}$ . \[ -a^3 + \left(2x-\frac{1}{2}\right)a^2+(-4x^2+4x)a-2x^2=0 \] Remember our original goal. It was to find an $x$ value such that $a=\frac{1}{2}$ is the only valid solution. Therefore, we can actually plug in $a=\frac{1}{2}$ back into the equation to look for values of $x$ such that the relation is satisfied, then eliminate undesirable answers. \[ 16x^2-10x+1=0 \] This is easily factored, allowing us to determine that $x=\frac{1}{8},\frac{1}{2}$ . The latter root is not our answer, since on line $\overline{AB}$ , $y\left(\frac{1}{2}\right)=0$ , the horizontal line segment running from $(0,0)$ to $(1,0)$ covers that point. From this, we see that $x=\frac{1}{8}$ is the only possible candidate. Going back to line $\overline{AB}, y= -\sqrt{3}x+\frac{\sqrt{3}}{2}$ , plugging in $x=\frac{1}{8}$ yields $y=\frac{3\sqrt{3}}{8}$ . The distance from the origin is then given by $\sqrt{\frac{1}{8^2}+\left(\frac{3\sqrt{3}}{8}\right)^2} =\sqrt{\frac{7}{16}}$ . That number squared is $\frac{7}{16}$ , so the answer is $\boxed{023}$ . ~Installhelp_hex Why this works: Assume the quartic equation has exactly two real roots. Then for $a=\frac{1}{2}$ to be the only solution, it must be a doubled root. Thus, even after dividing the quartic by $a-\frac{1}{2}$ , $a=\frac{1}{2}$ is still a root.
~inaccessibles Other explanation: Clearly $a=\frac{1}{2}$ works because then PQ = AB. Thus, we can factor out $a-\frac{1}{2}$ . From the 3rd degree polynomial, it must either have 3 real roots, or 1 real root and 2 non real roots, because it has real coefficients. So it must have atleast 1 real root. If this root wasn't $\frac{1}{2}$ this would mean their is another line PQ that satisfies the equation. Thus, their must be another root at $a=\frac{1}{2}$ and plugging in, we can solve for $x$ .
~Bigbrain_2009

\end{document}
